% O modelo inteiro em um pregão (22/02/2021), com os números e as séries que o
% próprio robô usou. Dados: dados/resultados/modelo_completo_2026-11-06/
% (ferramentas/ModeloCompleto.mq5, conferido contra o traço em conferencia.txt).
\documentclass[border=6pt]{standalone}
% Preâmbulo comum das figuras de ia/relatorios/modelo-completo e garch-funcionando
% (06/11/2026). Nenhum número é digitado aqui: vêm de valores.tex (\vX para o
% texto, \pX para coordenada) e dos CSVs de
% dados/resultados/modelo_completo_2026-11-06/ (ferramentas/ModeloCompleto.mq5)
% e dados/resultados/robo_em_acao_2026-10-30/ (traço do robô).
\usepackage[T1]{fontenc}
\usepackage{amsmath}
\usepackage{mathptmx}
\usepackage{pgfplots}
\usepgfplotslibrary{dateplot,fillbetween}
\usetikzlibrary{arrows.meta,positioning,calc,decorations.pathreplacing}
\pgfplotsset{compat=1.18}
\definecolor{cmod}{RGB}{0,90,181}     % ARIMA, ajuste
\definecolor{cdes}{RGB}{220,50,32}    % GARCH acima do limiar, veto
\definecolor{cdad}{RGB}{120,120,120}  % dado observado
\definecolor{cok}{RGB}{0,140,70}      % compra liberada
\definecolor{cgar}{RGB}{230,140,0}    % GARCH, limiar
\newcommand{\arqm}[1]{../../../dados/resultados/modelo_completo_2026-11-06/#1}
\newcommand{\arqr}[1]{../../../dados/resultados/robo_em_acao_2026-10-30/#1}
\input{\arqm{valores.tex}}
\input{\arqr{valores.tex}}
\newcommand{\passo}[1]{\tikz[baseline=(n.base)]\node[circle,fill=black!75,text=white,inner sep=1pt,font=\footnotesize\bfseries](n){#1};}
\pgfplotsset{
  mini/.style={table/col sep=comma, scale only axis, unbounded coords=jump,
    tick label style={font=\scriptsize}, label style={font=\scriptsize},
    grid=major, grid style={gray!20}, every axis plot/.append style={line width=0.7pt},
    /pgf/number format/use comma, /pgf/number format/1000 sep={.}},
  datas/.style={date coordinates in=x},
}
\tikzset{
  caixa/.style={draw=black!45, rounded corners=3pt, line width=0.6pt},
  titulo/.style={anchor=north west, font=\small\bfseries, inner sep=0pt},
  texto/.style={anchor=north west, font=\footnotesize, align=left, inner sep=0pt},
  seta/.style={-{Stealth[length=2.6mm]}, line width=1pt, black!55},
}

\begin{document}
\begin{tikzpicture}[x=1cm,y=1cm]
% ---------------------------------------------------------------- título
\node[anchor=north, font=\large\bfseries] at (8.5,0) {O robô ARIMA-GARCH em um pregão: \vSegDia};
\node[anchor=north, font=\small, text=black!70] at (8.5,-0.55)
  {decisão tomada na abertura, só com fechamentos já conhecidos, na configuração ARIMA-GARCH Sharpe};

% ---------------------------------------------------------------- 1 dados
\draw[caixa, fill=black!2] (0,-1.2) rectangle (17,-6.0);
\node[titulo] at (0.3,-1.4) {\passo{1} Os dados: a janela de 240 fechamentos que termina na véspera (refeita todo pregão)};
\begin{axis}[mini, datas, at={(0.9cm,-2.25cm)}, anchor=north west, width=6.6cm, height=2.3cm,
  xmin=2020-03-01, xmax=2021-02-28, xtick={2020-03-01,2020-09-01,2021-02-01},
  xticklabels={mar/2020,set/2020,fev/2021}, ymin=10, ymax=32,
  title={fechamento (R\$)}, title style={font=\scriptsize, yshift=-1.5mm}]
\addplot[cdad, line width=0.8pt] table[x=data, y=fechamento]{\arqm{janela_Seg.csv}};
\end{axis}
\draw[seta] (7.85,-3.4) -- node[above, font=\scriptsize, align=center]{log e\\diferença} (9.2,-3.4);
\begin{axis}[mini, datas, at={(9.6cm,-2.25cm)}, anchor=north west, width=6.9cm, height=2.3cm,
  xmin=2020-03-01, xmax=2021-02-28, xtick={2020-03-01,2020-09-01,2021-02-01},
  xticklabels={mar/2020,set/2020,fev/2021}, ymin=-40, ymax=25,
  title={log-retorno (\% ao dia)}, title style={font=\scriptsize, yshift=-1.5mm}]
\addplot[ycomb, cdad, line width=0.5pt] table[x=data, y=retorno]{\arqm{janela_Seg.csv}};
\end{axis}
\node[texto, text width=16.2cm] at (0.4,-5.25) {Fechamentos de \vSegIni{} a \vSegFim. A janela começa na queda de
  março de 2020 e termina na queda de 19/02/2021 (último fechamento, R\$~\vSegFech). O modelo trabalha com o log do preço e, como
  o passo~2 escolheu $d=1$, com a diferença dele, que é o log-retorno.};
\draw[seta] (8.5,-6.0) -- (8.5,-6.75);

% ---------------------------------------------------------------- 2 e 3 escolha do modelo
\draw[caixa, fill=black!2] (0,-6.8) rectangle (8.3,-11.55);
\draw[caixa, fill=black!2] (8.7,-6.8) rectangle (17,-11.55);
\node[titulo] at (0.3,-7.0) {\passo{2} Quantas diferenças? Teste ADF};
\node[titulo] at (9.0,-7.0) {\passo{3} Quantos termos? Grade do AIC};
% régua do ADF: v em [-7,0] -> x = 7.6 + v
\fill[cok!12] (0.6,-8.85) rectangle ({7.6+\pSegADFvc},-9.45);
\draw[black!60] (0.6,-9.45) -- (7.6,-9.45);
\foreach \v in {-7,-6,...,0} \draw[black!60] ({7.6+\v},-9.45) -- ({7.6+\v},-9.55) node[below, font=\tiny]{$\v$};
\draw[cdes, dashed, line width=0.8pt] ({7.6+\pSegADFvc},-8.0) -- ({7.6+\pSegADFvc},-9.45);
\node[font=\scriptsize, text=cdes, anchor=south] at ({7.6+\pSegADFvc},-8.05) {valor crítico \vSegADFvc};
\node[font=\scriptsize, text=cok!80!black, anchor=north west] at (0.65,-8.9) {estacionária};
\fill[cdad] ({7.6+\pSegADFnivel},-9.15) circle (2.6pt);
\node[font=\scriptsize, anchor=west] at ({7.75+\pSegADFnivel},-9.15) {log-preço \vSegADFnivel};
\fill[cmod] ({7.6+\pSegADFret},-8.6) circle (2.6pt);
\node[font=\scriptsize, anchor=west] at ({7.75+\pSegADFret},-8.6) {log-retorno \vSegADFret};
\node[texto, text width=7.7cm] at (0.3,-10.05) {O log-preço não passa no teste (anda sem rumo). O log-retorno
  passa. Então $d=1$: o modelo prevê o retorno, e não o preço.};
% grade do AIC
\begin{axis}[at={(9.55cm,-7.65cm)}, anchor=north west, scale only axis, width=2.75cm, height=2.75cm,
  xmin=-0.5, xmax=4.5, ymin=-0.5, ymax=4.5, xtick={0,...,4}, ytick={0,...,4},
  tick label style={font=\scriptsize}, label style={font=\scriptsize},
  xlabel={$q$ (erros)}, ylabel={$p$ (retornos)}, xlabel style={yshift=1.5mm}, ylabel style={yshift=-2mm},
  colormap={azul}{rgb255=(0,90,181) rgb255=(150,190,230) rgb255=(245,245,245)},
  point meta min=0, point meta max=20, enlargelimits=false, axis on top]
\addplot[matrix plot*, mesh/cols=5, point meta=explicit]
  table[col sep=comma, x=q, y=p, meta expr={min(\thisrow{delta},20)}]{\arqm{aic_Seg.csv}};
\draw[cdes, line width=1.4pt] (axis cs:{\pSegQ-0.5},{\pSegP-0.5}) rectangle (axis cs:{\pSegQ+0.5},{\pSegP+0.5});
\end{axis}
\node[texto, text width=4.15cm] at (12.75,-7.75) {25 modelos, $p$ e $q$ de 0 a~4. Quanto mais escuro, menor
  o AIC. O menor (em vermelho) foi o ARIMA(\vSegP,\vSegD,\vSegQ): os três últimos retornos e nenhum termo de erro.};
\node[anchor=north, font=\scriptsize\itshape, text=black!70, align=center, text width=16.5cm] at (8.5,-11.05)
  {passos 2 e 3: a cada 21 pregões. A última vez foi em \vSegSel, com os fechamentos de \vSegSelIni{} a \vSegSelFim};
\draw[seta] (4.15,-11.55) -- (4.15,-12.3);
\draw[seta] (12.85,-11.55) -- (12.85,-12.3);
\begin{scope}[yshift=-0.35cm]

% ---------------------------------------------------------------- 4 a 7: as duas pistas
\node[anchor=south, font=\small\bfseries, text=cmod] at (2.2,-11.75) {ARIMA $\Rightarrow$ direção};
\node[anchor=south, font=\small\bfseries, text=cgar!85!black] at (14.9,-11.75) {GARCH $\Rightarrow$ risco};
\draw[caixa, fill=cmod!5, draw=cmod!60] (0,-12.0) rectangle (8.3,-14.4);
\draw[caixa, fill=cgar!6, draw=cgar!70] (8.7,-12.0) rectangle (17,-14.4);
\node[titulo] at (0.3,-12.2) {\passo{4} Coeficientes (Hannan-Rissanen, todo pregão)};
\node[texto, text width=7.8cm] at (0.3,-12.75) {$\hat r_t = c + \phi_1 r_{t-1} + \phi_2 r_{t-2} + \phi_3 r_{t-3}$\\[3pt]
  $c=\vSegConst\%$, \ $\phi_1=\vSegPhiUm$, \ $\phi_2=\vSegPhiDois$, \ $\phi_3=\vSegPhiTres$\\[2pt]
  pesos estimados dos dados, um por retorno passado};
\node[titulo] at (9.0,-12.2) {\passo{6} GARCH(1,1) nos resíduos (todo pregão)};
\node[texto, text width=7.8cm] at (9.0,-12.75) {$\sigma^2_t = \omega + \alpha\,\varepsilon^2_{t-1} + \beta\,\sigma^2_{t-1}$\\[3pt]
  $\alpha=\vSegAlfa$ (choque de ontem), \ $\beta=\vSegBeta$ (memória)\\[2pt]
  $\varepsilon$ = o que o ARIMA errou em cada dia da janela};
\draw[seta, cgar!80!black] (8.3,-13.2) -- node[above, font=\scriptsize, text=black]{$\varepsilon$} (8.7,-13.2);
\draw[seta, cmod!70] (4.15,-14.4) -- (4.15,-14.95);
\draw[seta, cgar!80!black] (12.85,-14.4) -- (12.85,-14.95);

\draw[caixa, fill=cmod!5, draw=cmod!60] (0,-15.0) rectangle (8.3,-19.6);
\draw[caixa, fill=cgar!6, draw=cgar!70] (8.7,-15.0) rectangle (17,-19.6);
\node[titulo] at (0.3,-15.2) {\passo{5} Previsão do retorno de \vSegDia};
\node[titulo] at (9.0,-15.2) {\passo{7} Volatilidade prevista contra o limiar};
\begin{axis}[mini, datas, at={(0.95cm,-15.85cm)}, anchor=north west, width=6.95cm, height=2.35cm,
  xmin=2021-01-04, xmax=2021-02-24, xtick={2021-01-04,2021-02-01,2021-02-22},
  xticklabels={04/01,01/02,22/02}, ymin=-9, ymax=6, ylabel={\% ao dia}, ylabel style={yshift=-1mm}]
\addplot[ycomb, cdad, line width=1.6pt] table[x=data, y=retorno]{\arqm{janela_Seg.csv}};
\addplot[cmod, line width=0.8pt] table[x=data, y=ajustado]{\arqm{janela_Seg.csv}};
\addplot[only marks, mark=*, mark size=2.6pt, cmod] coordinates {(2021-02-22,\pSegPrev)};
\end{axis}
\node[texto, text width=7.8cm] at (0.3,-18.55) {Barras: retorno realizado. Linha: o que o modelo previa para cada dia.
  Ponto: previsão de \textbf{alta de \vSegPrev\%} $\Rightarrow$ sinal de compra.};
\begin{axis}[mini, datas, at={(9.65cm,-15.85cm)}, anchor=north west, width=6.95cm, height=2.35cm,
  xmin=2021-01-04, xmax=2021-02-24, xtick={2021-01-04,2021-02-01,2021-02-22},
  xticklabels={04/01,01/02,22/02}, ymin=0, ymax=8, ylabel={\% ao dia}, ylabel style={yshift=-1mm}]
\addplot[ycomb, cdad!70, line width=1.6pt] table[x=data, y expr={abs(\thisrow{residuo})}]{\arqm{janela_Seg.csv}};
\addplot[cgar!90!black, line width=0.9pt] table[x=data, y=sigma]{\arqm{janela_Seg.csv}};
\addplot[cdes, dashed, line width=0.8pt] coordinates {(2021-01-04,\pSegLim) (2021-02-24,\pSegLim)};
\addplot[only marks, mark=*, mark size=2.6pt, cdes] coordinates {(2021-02-22,\pSegVol)};
\end{axis}
\node[texto, text width=7.8cm] at (9.0,-18.55) {Barras: tamanho do erro. Linha: $\sigma$ do GARCH. Tracejado: limiar
  (percentil 70 da $\sigma$ na janela, \vSegLim\%). Ponto: \textbf{\vSegVol\% > \vSegLim\%} $\Rightarrow$ veto.};
\draw[seta, cmod!70] (4.15,-19.6) -- (4.15,-20.25);
\draw[seta, cgar!80!black] (12.85,-19.6) -- (12.85,-20.25);

% ---------------------------------------------------------------- 8 decisão
\draw[caixa, fill=black!2] (0,-20.3) rectangle (17,-24.35);
\node[titulo] at (0.3,-20.5) {\passo{8} Decisão (só compra, posição zerada no veto)};
\def\ya{-21.05}\def\yb{-21.6}\def\yc{-22.4}\def\yd{-23.2}
\def\xa{0.4}\def\xb{4.3}\def\xc{10.45}\def\xd{16.6}
\fill[cok!12] (\xb,\yb) rectangle (\xc,\yc);
\fill[cdes!15] (\xc,\yb) rectangle (\xd,\yc);
\fill[black!6] (\xb,\yc) rectangle (\xd,\yd);
\draw[black!40] (\xa,\ya) rectangle (\xd,\yd);
\draw[black!40] (\xa,\yb) -- (\xd,\yb) (\xa,\yc) -- (\xd,\yc) (\xb,\ya) -- (\xb,\yd) (\xc,\ya) -- (\xc,\yd);
\node[font=\footnotesize\bfseries] at ({(\xb+\xc)/2},{(\ya+\yb)/2}) {$\sigma$ prevista $\le$ limiar};
\node[font=\footnotesize\bfseries] at ({(\xc+\xd)/2},{(\ya+\yb)/2}) {$\sigma$ prevista $>$ limiar};
\node[font=\footnotesize\bfseries, text=cmod] at ({(\xa+\xb)/2},{(\yb+\yc)/2}) {previsão de alta};
\node[font=\footnotesize\bfseries, text=cmod] at ({(\xa+\xb)/2},{(\yc+\yd)/2}) {previsão de queda};
\node[font=\footnotesize] at ({(\xb+\xc)/2},{(\yb+\yc)/2}) {compra};
\node[font=\footnotesize, align=center] at ({(\xc+\xd)/2},{(\yb+\yc)/2}) {vetada: não compra e zera a posição};
\node[font=\footnotesize] at ({(\xb+\xd)/2},{(\yc+\yd)/2}) {zera a posição (o robô não vende a descoberto)};
\draw[cdes, line width=1.8pt] (\xc,\yb) rectangle (\xd,\yc);
\node[font=\scriptsize\bfseries, text=cdes, anchor=south east, fill=white, inner sep=1.5pt] at (\xd-0.05,\yb-0.05) {\vSegDia};
\node[texto, text width=16.2cm] at (0.4,-23.4) {O ARIMA diz para onde (linhas), o GARCH diz se é hora de estar exposto (colunas).
  Em \vSegDia{} o robô não comprou e zerou a compra feita na sexta. No pregão seguinte a janela anda um dia e
  tudo se refaz.};
% volta ao passo 1
\draw[seta, rounded corners=6pt, dashed] (0,-23.0) -- (-0.45,-23.0) -- node[sloped, above, font=\scriptsize]{próximo pregão} (-0.45,-3.25) -- (0,-3.25);
\end{scope}
\end{tikzpicture}
\end{document}
